\section{Selectivity no-go and its exact boundary}
\label{sec:boundary}

The force of \cref{thm:clone} comes from its information boundary, not from
the particular Boolean formula.  The following constructions remain inside
that boundary and are therefore clone-invariant:

\begin{itemize}
\item pair, triple, and arbitrary finite-arity join or lcm coherence;
\item interval-local and complete-source incidence M{\"o}bius transforms;
\item incidence-Hopf primitives, antipodes, and connected/logarithmic
  organization natural in the source;
\item partition-lattice cumulants after any source-natural moment has been
  fixed;
\item kernels that inspect the complete filtered tower rather than one
  interval;
\item the transported roof, metric, Gram, and compiled-operator stack; and
\item honest traces, determinant coefficients, or new functionals whenever
  they are constructed naturally from those same data.
\end{itemize}

Thus the theorem resolves the loophole left by a pair-local analysis.  A
hard-coded indicator of the complete integer-tower isomorphism class is also
covered: it takes value one on the isomorphic formal clone.  Declaring that
formal presentation to be “non-arithmetic” does not change the naturality
equation.

\begin{theorem}[Closure under formula changes]
\label{thm:formula-closure}
Within the category \(\Csrc\), no change of a natural join/M{\"o}bius kernel,
tuple arity, moment or cumulant convention, active-cutoff exhaustion,
finite-part convention, or same-object contraction can satisfy exact
integer-nonzero/every-UFD-zero selectivity.
\end{theorem}
\begin{proof}
Each change still defines an invariant natural in the same decorated object.
Apply \cref{thm:clone} to the integer source and its transported UFD clone,
then repeat the contradiction in \cref{cor:selectivity}.
\end{proof}

\begin{table}[t]
\centering
\small
\caption{Boundary of the clone theorem.  The right column lists new problems,
not exceptions silently available to SD-C32.}
\label{tab:boundary}
\begin{tabularx}{\textwidth}{Y Y}
\toprule
Covered by the theorem & Requires a new source object or weaker claim\\
\midrule
Any natural invariant of pointed divisibility, joins, intervals, M{\"o}bius
values, cutoffs, roof marks, Gram data, and admitted compiled coefficients &
A source-derived nonmultiplicative operation proved not to transport along
the valuation map\\
Local, nonlocal, nonlinear, complete-tower, and arbitrary finite-arity
constructions &
Addition--multiplication interaction, congruence correspondences, or a
genuine transfer operator with independent motivation\\
Any universally quantified class containing the transported formal UFD
clone &
Excluding the clone, thereby weakening “every UFD control”\\
Transported labels and coefficient marks &
Nontransportable printed-name or primality oracles, which violate the
information boundary\\
Exact baseline/control selectivity &
Abandoning exact selectivity or changing the target question\\
\bottomrule
\end{tabularx}
\end{table}

The most promising legitimate escape is therefore not another incidence
coefficient.  It is an independently defined operation relating two kinds
of source structure, followed by a proof that the formal UFD clone cannot
carry the same decoration by transport.  Such an operation would require a
new preregistration, control suite, analytic argument, and marker-ownership
audit.  It is reserved for a successor rather than inferred here.
